\documentclass[aspectratio=169]{beamer}
% Shared Beamer configuration for MUS2640 lecture decks (included after \documentclass).
\usepackage[utf8]{inputenc}
\usepackage[T1]{fontenc}
\usepackage{lmodern}
\usepackage{graphicx}
\usepackage{booktabs}

\usetheme{Madrid}
\usecolortheme{dolphin}
\usefonttheme{professionalfonts}

\setbeamercovered{transparent}
\setbeamertemplate{navigation symbols}{}
\setbeamertemplate{footline}[frame number]

\newcommand{\coursename}{MUS2640 · Sensing Sound and Music}
\newcommand{\instituteuo}{University of Oslo · Department of Musicology / RITMO}


\title{Week 6 · Time and rhythm}
\subtitle{\coursename}
\institute{\instituteuo}
\date{}

\begin{document}

\begin{frame}
  \titlepage
\end{frame}

\begin{frame}{Aims}
  \begin{itemize}
    \item Separate \textbf{pulse}, \textbf{meter}, \textbf{subdivision}, \textbf{groove}.
    \item Reason about micro-timing and syncopation across genres.
    \item Connect entrainment and dynamic attending to embodied listening.
  \end{itemize}
\end{frame}

\begin{frame}{Roadmap}
  \begin{enumerate}
    \item Periodic vs aperiodic time; notation vs performed timing.
    \item Meter as attentional behaviour; backbeats and idiomatic accents.
    \item Groove as pattern + performance approach; swing / microrhythm (examples).
    \item Entrainment: tapping, dance, ensemble coupling (hook to Body week).
  \end{enumerate}
\end{frame}

\begin{frame}{Embodied angle}
  \begin{itemize}
    \item Groove is often \textbf{felt} through movement — metre is not only intellectual.
    \item Small onset shifts carry style and affect (`in the pocket', laid-back\ldots).
  \end{itemize}
\end{frame}

\begin{frame}{Tools (optional)}
  \begin{itemize}
    \item Onset detection / tempograms — visual complements to ear (chapter demos).
    \item Critical listening: compare programmed vs performed swing feel.
  \end{itemize}
\end{frame}

\begin{frame}{In-class}
  \begin{itemize}
    \item Tap-along to excerpts; discuss where pulse sits vs snare/kick.
    \item Short reflection: does tapping change perceived groove?
  \end{itemize}
\end{frame}

\begin{frame}{Outside class}
  \begin{itemize}
    \item Listening analysis assignment slice; optional timing annotation exercise.
  \end{itemize}
\end{frame}

\begin{frame}{Summary}
  \begin{itemize}
    \item Rhythm research bridges signal processing, perception, and culture.
    \item Next: \textbf{Harmony and melody} — spectra, symbols, expectation.
  \end{itemize}
\end{frame}

\end{document}
